\section{Evaluation of Extractive KG Summaries}
\label{sec:benchmark}

Considering the magnitude of a KG, it is difficult---if not impossible---to manually extract a summary as the gold standard for evaluation. Instead, depending on the application, various \emph{quality metrics} have been used to quantitatively assess the quality of a summary from different perspectives. For example, Wang et al.~\shortcite{DBLP:journals/tkde/WangCPKQ23} devise an evaluation framework for measuring a summary's coverage of the classes, properties, EDPs, and LPs in the original KG. They also publish a benchmark called BANDAR with thousands of real-world KGs to be used for evaluation.
% For example, summaries extracted by maximizing pattern coverage are evaluated by the total proportion of covered classes and properties~\cite{DBLP:conf/wsdm/ChengJDXQ17,DBLP:journals/tweb/LiuCLQ19}, instantiated EDPs and LPs in the summary~\cite{DBLP:conf/semweb/00010LXPKQ21}, or similarity of class and property distribution between the summary and original KG~\cite{DBLP:conf/icde/SundaraAKDWCS10}.
Cheng et al.~\shortcite{DBLP:conf/wsdm/ChengJDXQ17} conduct a user study, inviting human users to rate and compare the quality of the summaries extracted by different approaches.

As to task-specific summaries, \emph{extrinsic evaluation} aims to assess the quality of a summary indirectly by testing its performance in downstream tasks using the summary. For example, Rietveld et al.~\shortcite{DBLP:conf/semweb/RietveldHSG14} and Guo et al.~\shortcite{DBLP:conf/service/GuoW21} measure a summary's coverage of query answers by comparing the results retrieved from the summary with the correct results retrieved from the original KG.
% For query-biased dynamic snippets, their query relevance is usually evaluated with the proportion of covered query keywords in the snippet~\cite{DBLP:conf/semweb/WangCK19,DBLP:conf/semweb/00010LXPKQ21}. Meanwhile, their compactness is measured with total edge weights~\cite{DBLP:conf/www/Shi0K20,DBLP:conf/semweb/ChengLZL20} or node weights~\cite{DBLP:conf/ijcai/Shi0T0K21,DBLP:conf/www/Shi0TKS21} of the summary. For personalized dynamic summary extraction, the metrics are usually associated with personalized importance measurement, such as the total importance of nodes contained in the result summary~\cite{DBLP:conf/esws/VassiliouAPK23}.
Wang et al.~\shortcite{DBLP:conf/semweb/00010LXPKQ21} conduct a user study to compare the usefulness of the summaries extracted by different approaches for assisting users in comprehending a large KG and then completing complex SPARQL queries over the KG. Useful summaries are expected to help complete SPARQL queries correctly and also quickly.

Besides, the \emph{time} for extracting a summary is an important factor in evaluating the practicability of an algorithm, especially for real-time applications such as KG search engines.
% Some works also report the storage cost for summaries and the compression rates compared to the original KG~\cite{DBLP:conf/semweb/00010LXPKQ21}.

% Existing approaches for extractive KG summarization typically focus on well-known open-sourced KGs with varied sizes and topics such as DBpedia\footnote{\url{https://www.dbpedia.org/}}, YAGO\footnote{\url{https://yago-knowledge.org/}} and WordNet\footnote{\url{https://wordnet.princeton.edu/}}. For summaries relying on keyword queries or query history for evaluation, existing works reuse structured query datasets~\cite{DBLP:conf/www/Shi0K20,DBLP:conf/ijcai/Shi0T0K21} or collect query logs from SPARQL endpoints~\cite{DBLP:conf/service/GuoW21}. To provide a general evaluation benchmark for extractive KG summaries, Wang et al.~\shortcite{DBLP:journals/tkde/WangCPKQ23} provide a dataset containing thousands of real-world datasets, as well as keyword queries derived from real-world data needs.